\documentclass[10pt]{article}
\usepackage[T1]{fontenc}
\usepackage[utf8]{inputenc}
\usepackage[french]{babel}
\usepackage{amsmath,amssymb,mathrsfs,mathpazo}
\usepackage{geometry}
\geometry{a4paper,margin=2cm}
\usepackage{graphicx,array,tabularx,tabularray,colortbl,multirow,multicol,booktabs,enumitem,pifont,ulem,fancybox,diagbox,cancel}
\usepackage{pstricks,pstricks-add,pst-plot,pst-tree,pst-node,pst-text,pst-func,pst-math,pst-bspline,pst-3dplot,pst-eucl,pst-solides3d}
\let\pscalculate\undefined
\usepackage{pst-bezier}
\usepackage{tikz,pgfplots,tkz-tab}
\usetikzlibrary{arrows,arrows.meta,patterns,calc,positioning,trees,decorations.markings,shapes.geometric,shapes.misc,shapes.arrows,intersections,decorations.pathmorphing,decorations.pathreplacing,backgrounds,shadows,babel}
\pgfplotsset{compat=1.18}
\usepgfplotslibrary{fillbetween}
\usepackage[np]{numprint}
\usepackage{listings,makecell,pas-tableur,twoopt}
\let\barre\undefined
\usepackage[pstricks]{tablvar}
\definecolor{darkgreen}{RGB}{0,100,0}
\providecommand{\red}{\color{red}}
\providecommand{\blue}{\color{blue}}
\providecommand{\green}{\color{green}}
\providecommand{\black}{\color{black}}
\providecommand{\white}{\color{white}}
\providecommand{\gray}{\color{gray}}
\providecommand{\cyan}{\color{cyan}}
\providecommand{\magenta}{\color{magenta}}
\providecommand{\yellow}{\color{yellow}}
\providecommand{\violet}{\color{violet}}
\providecommand{\orange}{\color{orange}}
\providecommand{\euro}{\texteuro}
\providecommand{\EUR}{\texteuro}
\providecommand{\ieme}{\textsuperscript{e}}
\providecommand{\degres}{\ensuremath{^\circ}}
\providecommand{\degree}{\ensuremath{^\circ}}
\providecommand{\vv}[1]{\overrightarrow{#1}}
\providecommand{\Coord}[1]{#1}
\providecommand{\decosix}{}
\providecommand{\starredbullet}{\ensuremath{\bullet}}
\pagestyle{empty}
\setlength{\parindent}{0pt}
\usepackage{adjustbox}\DeclareUnicodeCharacter{2236}{:}\DeclareUnicodeCharacter{2190}{\ensuremath{\leftarrow}}
\providecommand{\R}{\mathbb{R}}
\providecommand{\N}{\mathbb{N}}
\providecommand{\Z}{\mathbb{Z}}
\providecommand{\Q}{\mathbb{Q}}
\providecommand{\D}{\mathbb{D}}
\providecommand{\C}{\mathbb{C}}
\providecommand{\e}{{\mathrm{e}}}
\providecommand{\vect}[1]{\overrightarrow{#1}}
\providecommand{\vectt}[1]{\overrightarrow{\mathrm{#1}}}
\providecommand{\vvt}[1]{\overrightarrow{\mathrm{#1}}}
\providecommand{\np}[2][]{{#2~#1}}
\providecommand{\vv}[1]{{\overrightarrow{#1}}}
\renewcommand{\d}{{\mathrm{d}}}
\providecommand{\degres}{{^\circ}}
\providecommand{\degre}{{^\circ}}
\providecommand{\blue}{\color{blue}}
\providecommand{\red}{\color{red}}
\providecommand{\enskip}{\quad}
\providecommand{\llbracket}{\lbrack\!\lbrack}
\providecommand{\rrbracket}{\rbrack\!\rbrack}
\providecommand{\euro}{{\text{€}}}
\providecommand{\no}{{\text{n°}}}
\providecommand{\interval}[2]{{[#1;#2]}}
\providecommand{\intervFO}[2]{{[#1;#2[}}
\providecommand{\dotfill}{\ldots}

\providecommand{\ds}{\displaystyle}
\providecommand{\textding}[1]{#1}
\providecommand{\Oij}{\(\mathrm{O};\vec{\imath},\vec{\jmath}\)}
\providecommand{\Oijk}{\(\mathrm{O};\vec{\imath},\vec{\jmath},\vec{k}\)}
\providecommand{\Ouv}{\(\mathrm{O};\vec{u},\vec{v}\)}
\providecommand{\intervFF}[2]{\ensuremath{[#1;#2]}}
\renewcommand{\euro}{\ifmmode\text{\texteuro}\else\texteuro\fi}
\DeclareUnicodeCharacter{00B0}{\ifmmode^{\circ}\else\textdegree\fi}
\DeclareUnicodeCharacter{20AC}{\ifmmode\text{\texteuro}\else\texteuro\fi}
% Énoncés : Éducation nationale. Transcription et corrigés : APMEP, Denis Vergès et contributeurs. Découpage et classement : CoopMaths, Sébastien Lozano, Éric Elter et contributeurs. Adaptation de la présentation : Maths Explorer.
% CC BY-SA : https://coopmaths.fr/www/about/
% https://coopmaths.fr/alea/static/eam/2025/tex/eam_specifique_2025_06_sujet0v1_2.tex
\begin{document}
\noindent\textit{APMEP / CoopMaths. CC BY-SA. Présentation adaptée par Maths Explorer.}\par\medskip
\begin{minipage}[t]{0.65\textwidth}
Albert a acquis un étang d'une surface de 2\,000 m$^2$.

Le jour de son anniversaire, un dimanche, il installe des nénuphars sur une surface de 200 m$^2$.

\begin{enumerate}
\item Le dimanche d'après, la surface des nénuphars a augmenté de 40 m$^2$.
\begin{enumerate}
\item Quel pourcentage d'augmentation cela représente-t-il ?
\item Quelle est à présent la surface occupée par les nénuphars ?
\end{enumerate}
\end{enumerate}
\end{minipage}
\hfill
\begin{minipage}[t]{0.3\textwidth}
\vspace{0pt}
\begin{center}
\begin{tikzpicture}[scale=0.75]
\draw[fill=cyan!60, draw=black, line width=1pt] (0,0) ellipse (2.5cm and 1.5cm);
\draw[fill=green!70!black, draw=black, line width=1pt] (1.5,0) ellipse (1cm and 0.5cm);
\draw[-{Triangle[length=3mm,width=3mm]}, line width=1.25mm, black] (-1.5,1.6) -- (0,0.75);
\draw[-{Triangle[length=3mm,width=3mm]}, line width=1.25mm, black] (2.2,1.6) -- (1.5,0);
\node[above] at (-1.5,1.6) {Etang};
\node[above] at (2.2,1.6) {Nénuphars};
\end{tikzpicture}
\end{center}
\end{minipage}

\begin{enumerate}[resume]
\item Dans cette question, on suppose que la surface occupée par les nénuphars augmente de 40 m$^2$ chaque semaine, depuis la date de l'anniversaire, tant que cela est possible.
\begin{enumerate}
\item Quelle sera la surface occupée par les nénuphars 10 semaines après l'anniversaire ?
\item Est-il possible qu'un dimanche, la surface occupée par les nénuphars soit égale à 580 m$^2$ ? Justifier.
\item Au bout de combien de semaines, l'étang sera-t-il entièrement recouvert de nénuphars ?
\end{enumerate}

\item Dans cette question, on suppose que la surface occupée par les nénuphars augmente de 20\% chaque semaine, depuis la date de l'anniversaire, tant que cela est possible.
\begin{enumerate}
\item Quelle sera la surface occupée par les nénuphars 2 semaines après l'anniversaire ?
\item On considère un entier naturel $n$. Déterminer, en fonction de $n$, la surface occupée par les nénuphars $n$ semaines après l'anniversaire ?
\item Au bout de combien de semaines, l'étang sera-t-il entièrement recouvert par les nénuphars ? On pourra s'aider du tableau ci-dessous.

\begin{center}
\begin{tabular}{|c|c|c|c|c|c|c|c|c|c|}
\hline
$n = $ & 0 & 1 & 2 & 5 & 10 & 12 & 13 & 14 & 15 \\
\hline
$1,2^n \approx $ & 1 & 1,2 & 1,44 & 2,49 & 6,19 & 8,92 & 10,70 & 12,84 & 15,40 \\
\hline
\end{tabular}
\end{center}
\end{enumerate}

\item Réaliser sur votre copie un \textit{schéma} sur lequel apparaissent l'allure des nuages de points traduisant la progression de la surface occupée par les nénuphars, aussi bien dans le cas de la question 2 que dans le cas de la question 3, et faire figurer le moment où, dans chacun des cas, l'étang est recouvert par les nénuphars.
\end{enumerate}

\end{document}
